% BHU PYQ -- Transcribed & Corrected
% Paper: STAMN-21 : Probability Theory (Mid-Term)
% Exam: B.Sc. (Hons.) Semester II Mid-Term Examination, 2024-25 (NEP)
% Department: Department of Statistics

\documentclass[11pt,a4paper]{article}
\usepackage[a4paper,top=2.5cm,bottom=2.5cm,left=2.8cm,right=2.8cm,headheight=14pt]{geometry}
\usepackage{amsmath,amssymb}
\usepackage[T1]{fontenc}
\usepackage[utf8]{inputenc}
\usepackage{lmodern,microtype,fancyhdr,enumitem,hyperref,lastpage,parskip}

\hypersetup{colorlinks=true,urlcolor=black,linkcolor=black,pdftitle={STAMN21: Probability Theory Mid-Term}}
\pagestyle{fancy}\fancyhf{}\fancyhead[L]{\small\textit{BHU}}\fancyhead[R]{\small\textit{STAMN-21: Probability Theory}}\fancyfoot[C]{\small Page \thepage\ of \pageref{LastPage}}
\newcommand{\pts}[1]{\hfill\textit{[#1]}}

\begin{document}
\begin{center}
  {\large\bfseries BANARAS HINDU UNIVERSITY}\\[2pt]
  {\normalsize B.Sc. (Hons.) Semester II Mid-Term Examination, 2024-25 (NEP)}\\[6pt]
  {\large\bfseries DEPARTMENT OF STATISTICS}\\[4pt]
  {\large\bfseries Paper Code: STAMN-21 : Probability Theory (Mid-Term)}\\[4pt]
  {\normalsize Time: 1 Hour\hfill Maximum Marks: 20}
\end{center}
\rule{\linewidth}{0.4pt}\smallskip
\noindent\textit{Instructions: Attempt any FIVE questions. All questions carry equal marks (4 marks each).}
\bigskip

\noindent\textbf{Question 1.} State Axiomatic definition of Probability and prove $P(A \cup B) = P(A) + P(B) - P(A \cap B)$.\pts{4}

\medskip\rule{\linewidth}{0.2pt}\medskip

\noindent\textbf{Question 2.} Define conditional probability. State and prove Bayes' Theorem.\pts{4}

\medskip\rule{\linewidth}{0.2pt}\medskip

\noindent\textbf{Question 3.} Define Discrete Random Variable, PMF, and CDF. State properties of CDF.\pts{4}

\medskip\rule{\linewidth}{0.2pt}\medskip

\noindent\textbf{Question 4.} Define Continuous Random Variable and PDF. If $f(x) = k x^2 (1 - x)$ for $0 < x < 1$, find constant $k$.\pts{4}

\medskip\rule{\linewidth}{0.2pt}\medskip

\noindent\textbf{Question 5.} Define Mathematical Expectation $E(X)$ and Variance $V(X)$. Prove $V(aX + b) = a^2 V(X)$.\pts{4}

\vfill\rule{\linewidth}{0.4pt}
\begin{center}\small\href{https://bhu-exam.vercel.app/}{Click here to visit website}\end{center}
\end{document}
